\documentclass[11pt]{amsart}
\usepackage{geometry}                % See geometry.pdf to learn the layout options. There are lots.
\geometry{letterpaper}                   % ... or a4paper or a5paper or ... 
%\geometry{landscape}                % Activate for for rotated page geometry
%\usepackage[parfill]{parskip}    % Activate to begin paragraphs with an empty line rather than an indent
\usepackage{graphicx}
\usepackage{amssymb}
\usepackage{epstopdf}
\DeclareGraphicsRule{.tif}{png}{.png}{`convert #1 `dirname #1`/`basename #1 .tif`.png}

\usepackage{enumerate}
\usepackage{nicefrac}

\usepackage{mathrsfs}

\newcommand{\R}{\mathbb{R}}
\newcommand{\C}{\mathbb{C}}
\newcommand{\tstD}{\mathscr{D}}
\newcommand{\tstE}{\mathscr{E}}
\newcommand{\tstS}{\mathscr{S}}
\newcommand{\tstC}{\mathscr{C}}
\newcommand{\semP}{\mathscr{P}}
\newcommand{\eps}{\varepsilon}
\newcommand{\supp}{\mathrm{supp}}
\newcommand{\singsupp}{\mathrm{singsupp}}
\newcommand{\exd}{\mathrm{d}}

\title{Math 581 Assignment 5}
\author{Due Wednesday April 9}
\date{Winter 2014}                                           % Activate to display a given date or no date

\begin{document}
\maketitle

\begin{enumerate}[1.]
\item
For each of the following functions,
determine if it is a tempered distribution,
and if so compute its Fourier transform.
\begin{enumerate}[{\em a)}]
\item
$x\sin x$,
\item
$\frac1x\sin x$,
\item
$e^{i|x|^2}$, 
\item
$x\vartheta(x)$, where $\vartheta$ is the Heaviside step function,
\end{enumerate}
\item
For $t>0$, let $\varkappa_t:\R^n\to\R^n$ be the dilation operator $\varkappa_t(x)=tx$.
We call a function $f$ defined either on $\Omega=\R^n$ or on $\Omega=\R^n\setminus\{0\}$ {\em homogeneous of degree $s\in\C$} if 
$\varkappa_t^*f=t^sf$ for all $t>0$, that is, $f(tx)=t^sf(x)$ for all $t>0$ and $x\in\Omega$.
\begin{enumerate}[{\em a)}]
\item
Extend the dilation operator $\varkappa_t$ to distributions on $\Omega$, and give a natural definition of homogeneous distributions.
\item
Show that differentiation of distributions preserves homogeneity.
\item
Derive a formula for $\widehat{u\circ A}$, where $A$ is an $n\times n$ invertible matrix.
\item
Show that the Fourier transform of a homogeneous distribution of degree $s$ is
a homogeneous distribution of degree $-s-n$.
\item
From homogeneity considerations, compute the inverse Fourier transform of $|\xi|^{-2}$ up to a constant factor.
Then determine this constant factor.
Take the space dimension to be $n\geq3$.
\end{enumerate}
\item
\begin{enumerate}[{\em a)}]
\item
There are (at least) two ways to define the Fourier transform on $L^2(\R^n)$.
\begin{itemize}
\item
Extend the Fourier transform from $\tstS$ to $L^2$ by using the density of $\tstS$ in $L^2$ (as well as the Plancherel bound).
\item
First define the Fourier transform on $\tstS'$ by duality, and then restrict it to $L^2$.
\end{itemize}
Show that these two approaches are consistent with each other.
\item
Show that the Fourier transform acting on $L^1$ is not onto $\mathscr{C}_0$.
\end{enumerate}
\item
\begin{enumerate}[{\em a)}]
\item
Give an example of $u\in\tstC(\R^n)$ such that $\varphi\mapsto\int u\varphi$ is a tempered distribution and that there is no polynomial $p$ satisfying 
$|u(x)|\leq|p(x)|$ for all $x\in\R^n$.
\item
Show that the Radon measure
$$
\varphi\mapsto\sum_{k=1}^\infty a_k \varphi(k) ,
$$
is in $\tstS'(\R)$ if and only if there are constants $m\geq0$ and $c>0$ such that
$$
|a_k| \leq ck^m ,
\qquad \textrm{for all} \,\, k .
$$
\end{enumerate}
\item
\begin{enumerate}[a)]
\item
Prove that if $p$ is a polynomial with no real zeroes, then there are constants $c>0$ and $m$ such that $|p(\xi)|\geq c(1+|\xi|)^m$ for all $\xi\in\R^n$.
Operators $p(D)$ with $p$ satisfying this condition are called {\em strictly elliptic}.
\item
Show that if $p(D)$ strictly elliptic, then the equation $p(D)u=f$ has a solution for each $f\in\tstS'$.
\end{enumerate}
\item
Prove the following.
\begin{enumerate}[a)]
\item
The {\em Paley-Wiener-Schwartz theorem}:
Let $K\subset\R^n$ be a compact convex set, and let $\psi\in\tstS'$.
Then a necessary and sufficient condition for $\psi$ to be the Fourier transform of a distribution supported in $K$
is that $\psi$ is entire and satisfies the growth estimate
$$
|\psi(\zeta)| \leq C(1+|\zeta|)^{N} e^{I_K(\eta)},
\qquad \zeta=\xi+i\eta\in\C^n,
$$
with some constants $C$ and $N$.
Hence the Fourier-Laplace transform of a compactly supported distribution is an entire function of exponential type.
Recall that the indicator function $I_K$ of $K$ is defined as
$$
I_K(\eta) = \sup_{x\in K} \eta\cdot x.
$$
\item
If the set of real zeroes of $p$ is bounded, then every tempered distribution solution of $p(D)u=0$ is an entire function of exponential type.
\item
If $p(D)$ is a nontrivial differential operator and $u\in\tstE'$ satisfies $p(D)u=0$ then $u=0$.
\end{enumerate}
\item
Let $p$ be a nonzero polynomial. Show the following.
\begin{enumerate}[a)]
\item
The equation $p(D)u=f$ has at least one smooth solution for every $f\in\tstD$.
\item
If all solutions of $p(D)u=0$ are smooth, then $p(D)$ is hypoelliptic.
\item
If $p(D)$ admits a fundamental solution that is smooth outside some ball of finite radius (centred at the origin), 
then $p(D)$ is hypoelliptic.
\end{enumerate}
\item
Recall that by H\"ormander's theorem, $p(D)$ is hypoelliptic if and only if 
for any $\eta\in\R^n$ one has $p(\xi+i\eta)\neq0$ for all sufficiently large $\xi\in\R$.
\begin{enumerate}[a)]
\item
Construct a non-hypoelliptic polynomial $p$ in dimension $n>1$ such that
$|p(\xi)|\to\infty$ as $|\xi|\to\infty$ for $\xi\in\R^n$.
\item
For any given $c>0$, construct a non-hypoelliptic polynomial $p$ in dimension $n>1$ such that
$|p(\xi+i\eta)|\to\infty$ uniformly in $\{|\eta|\leq c\}$ as $|\xi|\to\infty$ for $\xi\in\R^n$.
\end{enumerate}
\end{enumerate}

\end{document}  












